
\documentclass{article}
\usepackage[utf8]{inputenc}
\usepackage{graphicx}
\usepackage{hyperref}
\usepackage{amsmath}
\usepackage{geometry}
\geometry{a4paper, margin=1in}

\title{Multiplicación de matrices de forma secuencial y paralela}
\author{Presentado por: \\ Estudiante}
\date{2022}

\begin{document}

\maketitle

\begin{abstract}
Por medio del presente documento se muestra el proceso realizado para
llevar a cabo el algoritmo de multiplicación de matrices en el lenguaje de
programación C, las multiplicaciones fueron realizadas de dos formas:
secuencial y paralela, esto con el fin de analizar su desempeño tras
ejecutarlo, variando el tamaño de las matrices a multiplicar y la disposición
de los recursos del procesador a través de la librería pthread. Se han
realizado pruebas ejecutando el algoritmo de forma secuencial, luego se
implementaron hilos con la librería pthread y 32 hilos y por último se
implementó procesos con forks. Con base en estas pruebas se han obtenido
las respectivas conclusiones
\end{abstract}

\tableofcontents

\section{Introducción}
A lo largo de los años la computación ha ido evolucionando hasta el punto de poder
resolver problemas de grandes dimensiones en un tiempo más reducido, un ejemplo
de estos problemas es la multiplicación de matrices. A pesar de la brevedad para
programar este algoritmo en algún lenguaje de programación, “se ha invertido
mucho trabajo en hacer que los algoritmos de multiplicación de matrices sean
eficientes. [...] Se han diseñado muchos algoritmos diferentes para multiplicar
matrices en diferentes tipos de hardware, incluidos sistemas paralelos y distribuidos,
donde el trabajo computacional se distribuye en múltiples procesadores” (Qaz Wiki,
2020). En este caso se ha optado por resolver el problema programando dos
versiones, una para darle solución de manera secuencial y otra para darle solución
de manera paralela. Por otra parte, teniendo en cuenta la definición del speedup (o
aceleración) “Herramienta de la Mejora continua que de forma metodológica busca
incrementar la velocidad de procesamiento de las máquinas” (Progressa Lean,
2014) se pretende calcular este valor en base a los resultados arrojados tras la
ejecución de ambas versiones del algoritmo programadas en el lenguaje C,
utilizando los conceptos vistos durante las sesiones del curso HPC.

Finalmente se busca comprobar que, ejecutando el algoritmo con múltiples hilos, se
hace posible reducir los tiempos de ejecución del algoritmo, beneficiando a las áreas
donde este se aplica, tales como el álgebra lineal o la inteligencia artificial.

\section{Marco conceptual}

\subsection{High Performance Computing}
Cuando hablamos de HPC por sus siglas en inglés, o Computación de alto
rendimiento en español, “hacemos referencia a un campo de la computación
actual que da solución a problemas tecnológicos muy complejos y que
involucran un gran volumen de cálculos o de coste computacional.” (López,
2017) Para nuestro caso, al aumentar la dimensión de las matrices a
multiplicar, esto supone un alto costo computacional, el cual, según el
análisis, se ve disminuido gracias al uso de herramientas que permiten,
precisamente, disminuir el tiempo de ejecución.

\subsection{Multiplicación matricial}
Para implementar el algoritmo en el lenguaje C es necesario comprender
previamente algunos conceptos, entre ellos, como es el proceso de
multiplicar matrices de forma manual, el cual podemos aprender fácilmente
gracias a vídeos explicativos y a diversos sitios web, tal como Producto de
matrices.
Otros sitios en internet nos dan una idea más general de la programación de
la multiplicación matricial, como por ejemplo el que se expone a
continuación, que ha sido extraído del blog Análisis y diseño de algoritmos.
Dadas dos matrices A y B, tales que el número de columnas de la matriz A
es igual al número de filas de la matriz B, es decir:
\[ A := (a_{ij})_{m \times n} \quad Y \quad B := (b_{ij})_{n \times p} \]

La multiplicación de A por B genera una nueva matriz resultado de la
forma: \( C = AB := (C_{ij})_{m \times p} \) donde cada elemento \( C_{ij} \) está definido por:
\[ C_{ij} = \sum_{k=1}^{q} a_{ik}b_{kj} \]

Definición que permite acercarse más a una idea del código ya que de
entrada plantea las tres variables a iterar para cada una de las celdas
de la matriz resultado.

\subsection{Complejidad Computacional}
“La Teoría de la Complejidad estudia la eficiencia de los algoritmos en
función de los recursos que utiliza en un sistema computacional (usualmente
espacio y tiempo)” (Cortéz, 2004). Para este estudio, la complejidad
computacional permite observar que en la multiplicación de matrices se
gastan bastantes recursos de procesamiento en cuanto a tiempo, es decir,
para matrices de dimensiones muy grandes, el procesador estará ocupado
bastante tiempo resolviendo la multiplicación, lo que motiva aún más a
programar de manera paralela el algoritmo, aprovechando al máximo los
recursos del procesador.

\subsection{Programación paralela}
"La computación paralela es el uso de múltiples recursos computacionales
para resolver un problema. Se distingue de la computación secuencial en
que varias operaciones pueden ocurrir simultáneamente.” (Ferestrepoca,
2019) En nuestro caso de estudio se ejecutará el código de multiplicación de
matrices de forma secuencial y de forma paralela, esto con el fin de
comparar los resultados obtenidos por cada ejecución y de esta manera
poder analizar la eficiencia de la programación paralela.

\subsection{Hilos}
"La diferencia fundamental entre hilos y procesos es que un proceso
dispone de un espacio de memoria separado, mientras que un hilo no. De
esta forma los hilos pueden trabajar directamente con las variables creadas
por el proceso padre, mientras que los procesos hijos no pueden hacerlo.”
(Montero \& Antunez, 2011) Esta información nos facilita la comprensión del
principio aplicado para esta implementación.

\subsection{Speedup}
El speedup representa la ganancia que se obtiene en la versión paralela de
un programa con respecto a su versión secuencial. Para este estudio se ha
tomado el speedup como:
\[ \text{Speedup} = \frac{T^*(n)}{T_p(N)} \]
Donde \( T^*(n) \) hace referencia al tiempo de ejecución secuencial
y \( T_p(N) \) hace referencia al tiempo de ejecución paralelo.

\section{Marco Contextual}

\subsection{Características de la máquina}
Se probó el algoritmo para la multiplicación de matrices generando matrices
cuadradas de dimensiones 10, 100, 200, 400, 800, 1600 y 3200
respectivamente. Uno de los algoritmos fue ejecutado de manera secuencial
y el otro utilizando hilos. Se analizaron los distintos métodos evaluando los
recursos consumidos como el tiempo; se sacaron algunas gráficas para
mostrar los resultados obtenidos.

Características del PC donde se realizaron las pruebas:
\begin{itemize}
    \item Procesador: Intel(R) Core(TM) i7-10700
    \item Cantidad de núcleos: 8
    \item Cantidad de subprocesos/hilos: 16
    \item Frecuencia básica del procesador: 2.90GHz
    \item Frecuencia Turbo máxima: 4.80GHz
    \item Ram 16GB DDR4 @ 2400 MHz
    \item SSD: 240GB - R: 540 MB/s - W: 500 MB/s
    \item Sistema operativo: Ubuntu 20.04.4 LTS
\end{itemize}

\subsection{Desarrollo}
Para iniciar el desarrollo de esta actividad lo primero era generar matrices
cuadradas con números random, por lo cual se usó la librería stdlib.h. El
tamaño de las matrices cuadradas usadas fueron 10, 100, 200, 400, 800,
1600 y 3200. Con las matrices a multiplicar listas, se pudo aprovechar la
implementación en forma secuencial, con hilos y por último con procesos
para analizar el comportamiento del algoritmo.
El código en C de la multiplicación de matrices utilizando hilos se basa en el
código \href{https://github.com/ferestrepoca/HPC-1/blob/master/Teorico/Multithread/pthreads.c}{Programación multithread en C (pthread.h)} con el cual se comenzó
añadiendo las librerías de hilos pthread.h para crear, asignar y juntar las
actividades que debe realizar cada hilo y la librería time.h para poder
contabilizar el tiempo que demora la ejecución. Para trabajar con time.h de
forma correcta, algunas partes del código se basan en el recurso \href{https://www.geeksforgeeks.org/how-to-measure-time-taken-by-a-program-in-c/}{8 formas de medir el tiempo de ejecución en C / C ++}. El código empieza asignando
el número de hilos y las variables que almacenarán el tamaño de las filas y
las columnas de las matrices que se usarán en la ejecución de la
multiplicación. Se pide al usuario que digite el tamaño de la matriz para
seguido a esto, con la función malloc reservar memoria para poder trabajar
con hilos, posteriormente se procede a la lectura de las matrices y la
creación de los hilos a través de funciones implementadas dentro del
código. Una vez creados los hilos, se toma el tiempo a través de
gettimeofday. Por otra parte, cuando ya está definida la forma en cómo se
creará la cantidad de hilos y la unión del trabajo de cada uno, se hace la
declaración de la función multiplicación la cual está diseñada para asignar
las filas por multiplicar en cada uno de los hilos que la utilicen, además,
dentro de esta función se realiza la operación de multiplicación de las
matrices por medio de tres ciclos anidados. Finalmente, se escribe el
resultado de la multiplicación de matrices hecho por cada hilo, se une a
través de la función proporcionada por pthread.h para tener la matriz
resultante completa y utilizando nuevamente la función gettimeofday se
hace posible calcular el tiempo de ejecución entre la creación de hilos y la
unión de sus resultados. La salida de este programa es únicamente el
tiempo de ejecución. El código de la multiplicación de matrices secuencial
usa las mismas librerías usadas con los hilos, excepto pthread.h, ya que
esta vez la ejecución es secuencial. El código empieza definiendo las
variables que almacenarán el tamaño de filas y columnas de la matriz que
vamos a usar en la ejecución de la multiplicación, se declaran los arreglos
bidimensionales para las matrices a multiplicar siendo estas DoubleA y
DoubleB y el arreglo bidimensional Double C para el resultado. Se le pide al
usuario que ingrese el tamaño de la matriz y se usa la función malloc para
hacer la reserva de memoria para poder ejecutar y se procede a la lectura
de las matrices a través de la función read mat, la cual recibe la cantidad de
filas y columnas de la matriz. Posteriormente se toma el tiempo con
gettimeofday y se hace el llamado a la función matmul cuyos parámetros de
entrada son las filas y columnas, que representan la dimensión de las
matrices a multiplicar y dentro de esta función se realiza la operación de
multiplicación de las matrices por medio de tres ciclos anidados y una vez
termina esta función de ejecutarse, se toma el tiempo nuevamente con
gettimeofday, con el fin de establecer la diferencia de tiempo entre el inicio y
fin de la ejecución. Al igual que en el caso de los hilos, este programa
secuencial imprime únicamente el tiempo. Luego de hacer la ejecución de
cada multiplicación de matriz con diferentes tamaños y en el caso de los
hilos con diferente número de hilos, obtenemos los resultados del tiempo de
ejecución de cada uno de estos, lo que permite realizar los análisis y
cálculos de la variable speedup para cada uno de los casos.
El código fuente resultante de la implementación descrita se encuentra
registrado en el enlace: \href{https://github.com/ferestrepoca/HPC-1/tree/master/Practico/Caso%20de%20Estudio%20Matrices}{Caso de Estudio Matrices} con las respectivas
instrucciones para una ejecución correcta.

\section{Pruebas}
Ahora, para complementar el análisis de los datos y tener una vista un poco
más amplia, se han utilizado las tablas para hacer gráficas comparativas que
nos permitan visualizar la eficiencia de trabajar con hilos.

\begin{figure}[h!]
\centering
\includegraphics[width=\textwidth]{graficas/tiempo_ejecucion.png}
\caption{Tiempo de ejecución en función de las dimensiones de la matriz.}
\end{figure}

A continuación se muestra una gráfica comparativa entre todos los
speed up de cada método de paralelización usado, es bueno recordar que el
speedup es la división de el tiempo que se tiene en la ejecución secuencial,
sobre el tiempo obtenido al paralelizar, por lo que en la gráfica no se incluye
en speed up del método secuencial ya que este es una constante igual a 1.

\begin{figure}[h!]
\centering
\includegraphics[width=\textwidth]{graficas/speedup.png}
\caption{Valor de Speed Up en función de las dimensiones de la matriz.}
\end{figure}

\section{Conclusiones}
\begin{enumerate}
    \item Luego de analizar los datos obtenidos, se puede concluir que el uso
    de hilos es muy necesario para trabajar con matrices de dimensiones
    superiores a 1200, ya que a partir de este tamaño el tiempo de
    ejecución en secuencial es mucho mayor.
    \item Los valores de speedup de las ejecuciones con hilos, aumentan
    proporcionalmente al aumento de los hilos utilizados
    \item Usar más de ocho hilos es innecesario, debido a que la diferencia
    entre los speedup obtenidos con las ejecuciones de 16 y 32 hilos, es
    muy pequeña.
    \item Usar procesos no es tan recomendable, ya que el speed up que se
    obtiene con pocas o altas dimensiones, no sobresale demasiado
    frente a paralizaciones superiores a 4 hilos
    \item Analizando la gráfica 1, se nota un pico cerca a las ejecuciones con
    matrices de dimensiones 800, esto, según suposición nuestra, debido
    al proceso de asignación de memoria que es necesario hacer para
    poder trabajar con matrices de dimensión superior a 830
    aproximadamente
\end{enumerate}

\section{Bibliografía}
\begin{thebibliography}{9}
    \bibitem{cortez}
    Cortéz, A. (2004). TEORÍA DE LA COMPLEJIDAD COMPUTACIONAL Y TEORÍA DE.
    \textit{Revista De Investigación De Sistemas E Informática}, 1(1), 102-105. Recuperado de:
    \url{https://revistasinvestigacion.unmsm.edu.pe/index.php/sistem/article/view/3216}

    \bibitem{ferestrepoca}
    Ferestrepoca. (2019). Recuperado de: \url{http://ferestrepoca.github.io/paradigmas-de-programacion/paralela/paralela_teoria/index.html}

    \bibitem{lopez}
    López, A. d. (2017, 11 30). \textit{Encamina}. From Por una nube sostenible:
    \url{https://blogs.encamina.com/por-una-nube-sostenible/introduccion-al-high-performance-computing/}

    \bibitem{montero}
    Montero, L. H., \& Antunez, R. R. (2011, 07). \textit{Parallel programming: definitions, mechanisms
    and trouble}. From
    \url{https://www.researchgate.net/publication/274960405_Parallel_programming_definitions_mechanisms_and_trouble}

    \bibitem{progressa}
    Progressa Lean. (2014, 11 2). \textit{Progressa Lean}. From \url{https://www.progressalean.com/que-es-speed-up/}

    \bibitem{qazwiki}
    Qaz Wiki. (2020). \textit{Qaz Wiki}. From \url{https://es.qaz.wiki/wiki/Matrix_multiplication_algorithm}
\end{thebibliography}

\end{document}
